\section{Evidence from the recorded source}\label{sec:evidence}

One script in the paper's source tree produces every result in this section.
It runs in float64 on CPU with fixed seeds, against STATADIC
\jcvalue{statadic-version} with package source \jcvalue{package-source-state}
relative to base revision \texttt{\jcvalue{revision}}, JAX
\jcvalue{jax-version} and
Python \jcvalue{python-version}. It also executes \cref{lst:ols,lst:poisson}
verbatim.

\subsection{Numerical agreement with independent software}\label{sec:agreement}

\Cref{tab:agreement} compares STATADIC with statsmodels
\jcvalue{statsmodels-version}, dcor \jcvalue{dcor-version} and lifelines
\jcvalue{lifelines-version} on synthetic fixtures:
\begin{itemize}
  \item a heteroskedastic linear model with $n = 400$;
  \item a Poisson model with exposure and $n = 400$;
  \item two Gaussian samples of sizes 200 and 150, for energy distance;
  \item a Weibull accelerated failure time (AFT) model with $n = 400$ and
    \jcvalue{censored-share}\% right censoring.
\end{itemize}

Differences are at rounding level except for the survival model, where each
library stops an iterative optimizer at its own tolerance. One convention
differs: the HC1 standard errors of the statsmodels GLM equal its HC0 values,
whereas STATADIC's HC1 applies the factor $n/(n-p)$ to covariance, and hence
$\sqrt{n/(n-p)}$ to standard errors. Without that factor the HC1 difference is
\jcnum{hc1-unscaled-difference}. Agreement on a fixture is
evidence of numerical conformance for that configuration, not of statistical
validity.

\begin{table}[t]
  \centering
  \small
  \caption{Maximum absolute differences between STATADIC and independent software
  in float64.}
  \label{tab:agreement}
  \setlength{\tabcolsep}{4pt}
  % AUTO-GENERATED by analysis/generate.py; do not edit.
\begin{tabular}{@{}llrr@{}}
\toprule
Quantity & Comparator & Values & Max.\ abs.\ difference \\
\midrule
OLS coefficients & statsmodels 0.14.6 & 2 & \num{1.1e-16} \\
OLS standard errors: classical, HC0--HC3 & statsmodels 0.14.6 & 10 & \num{4.2e-17} \\
Poisson coefficients with exposure & statsmodels 0.14.6 & 2 & \num{3.3e-16} \\
Poisson log likelihood and deviance & statsmodels 0.14.6 & 2 & \num{1.1e-13} \\
Poisson standard errors: model-based, HC0 & statsmodels 0.14.6 & 4 & \num{3.3e-13} \\
Poisson HC1 standard errors\textsuperscript{a} & statsmodels 0.14.6 & 2 & \num{6.9e-14} \\
Poisson mean predictions under new exposure & statsmodels 0.14.6 & 3 & \num{2.7e-15} \\
Euclidean energy distance & dcor 0.7 & 1 & \num{0} \\
Weibull AFT parameters, right censoring & lifelines 0.30.0 & 3 & \num{7.1e-07} \\
Weibull AFT log likelihood & lifelines 0.30.0 & 1 & \num{2.7e-10} \\
\bottomrule
\end{tabular}


  \smallskip
  \footnotesize\textsuperscript{a}Compared after applying $\sqrt{n/(n-p)}$ to
  statsmodels' HC0 values; see the text.
\end{table}

\subsection{Covariance choice and coverage}\label{sec:coverage}

\Cref{tab:agreement} shows that STATADIC computes the HC estimators correctly. It
does not show that intervals built from them cover. To illustrate the
difference, we generated $y_i = 1 + 0.5 x_i + \sqrt{x_i}\,\varepsilon_i$, with
$\log x_i$ and $\varepsilon_i$ independent standard normal. Sample sizes ran
from 25 to 1600, with \jcvalue{trials} trials each. Each trial fits least
squares once and attaches HC0, HC1, HC2 and HC3 covariance to that same fit;
the slope interval uses the normal reference at 95\%, and an unavailable
interval counts as uncovered. The study uses STATADIC's numerical estimator bricks
under \texttt{jax.jit} and \texttt{jax.vmap}, so one compiled program evaluates
every trial of a cell.

\begin{figure}[t]
  \centering
  \pgfplotsset{every axis/.append style={ymin=\jcvalue{coverage-ymin},ymax=1,
      tick label style={font=\footnotesize},label style={font=\footnotesize},
  title style={font=\small}}}
  \begin{tikzpicture}
\begin{axis}[width=6.6cm,height=5cm,xlabel={Sample size},ylabel={Slope coverage},title={HC0},xtick={0,1,2,3,4,5,6},xticklabels={{25},{50},{100},{200},{400},{800},{1600}}]
% Interval meaning: Wilson 95% interval for coverage over 10000 trials
\addplot+[no marks,forget plot] coordinates {(0,0.6825768848970969) (0,0.7006758901275857)};
\addplot+[no marks,forget plot] coordinates {(1,0.7505189879210942) (1,0.7672821009228953)};
\addplot+[no marks,forget plot] coordinates {(2,0.7998580583286615) (2,0.815305629074024)};
\addplot+[no marks,forget plot] coordinates {(3,0.8443955232747883) (3,0.8583345258706262)};
\addplot+[no marks,forget plot] coordinates {(4,0.8719508884456331) (4,0.8847584247877647)};
\addplot+[no marks,forget plot] coordinates {(5,0.8944785166197612) (5,0.9062139006853268)};
\addplot+[no marks,forget plot] coordinates {(6,0.908447599852065) (6,0.9194343726973405)};
\addplot+[only marks,forget plot] coordinates {(0,0.6917) (1,0.759) (2,0.8077) (3,0.8515) (4,0.8785) (5,0.9005) (6,0.9141)};
\addplot[black,dashed,no marks,forget plot] coordinates {(0,0.95) (6,0.95)};
\end{axis}
\end{tikzpicture}
\hfill\begin{tikzpicture}
\begin{axis}[width=6.6cm,height=5cm,xlabel={Sample size},ylabel={Slope coverage},title={HC1},xtick={0,1,2,3,4,5,6},xticklabels={{25},{50},{100},{200},{400},{800},{1600}}]
% Interval meaning: Wilson 95% interval for coverage over 10000 trials
\addplot+[no marks,forget plot] coordinates {(0,0.6985055902989736) (0,0.7163350503771957)};
\addplot+[no marks,forget plot] coordinates {(1,0.7569937371114757) (1,0.7736024365533691)};
\addplot+[no marks,forget plot] coordinates {(2,0.8024970904401808) (2,0.8178646001709774)};
\addplot+[no marks,forget plot] coordinates {(3,0.8468413792283328) (3,0.8606868267249023)};
\addplot+[no marks,forget plot] coordinates {(4,0.8727687765377132) (4,0.8855399222982914)};
\addplot+[no marks,forget plot] coordinates {(5,0.8949914473662954) (5,0.906700585940422)};
\addplot+[no marks,forget plot] coordinates {(6,0.9085504560920603) (6,0.9195314396576713)};
\addplot+[only marks,forget plot] coordinates {(0,0.7075) (1,0.7654) (2,0.8103) (3,0.8539) (4,0.8793) (5,0.901) (6,0.9142)};
\addplot[black,dashed,no marks,forget plot] coordinates {(0,0.95) (6,0.95)};
\end{axis}
\end{tikzpicture}


  \medskip
  \begin{tikzpicture}
\begin{axis}[width=6.6cm,height=5cm,xlabel={Sample size},ylabel={Slope coverage},title={HC2},xtick={0,1,2,3,4,5,6},xticklabels={{25},{50},{100},{200},{400},{800},{1600}}]
% Interval meaning: Wilson 95% interval for coverage over 10000 trials
\addplot+[no marks,forget plot] coordinates {(0,0.7772446798874184) (0,0.7933361338425993)};
\addplot+[no marks,forget plot] coordinates {(1,0.8123476279170886) (1,0.8274066131256782)};
\addplot+[no marks,forget plot] coordinates {(2,0.846229855766208) (2,0.8600988109850718)};
\addplot+[no marks,forget plot] coordinates {(3,0.8767572927270054) (3,0.8893484109217079)};
\addplot+[no marks,forget plot] coordinates {(4,0.8935553559860696) (4,0.9053377525160855)};
\addplot+[no marks,forget plot] coordinates {(5,0.9099906788405526) (5,0.9208901417137408)};
\addplot+[no marks,forget plot] coordinates {(6,0.9177140205011988) (6,0.9281610400775347)};
\addplot+[only marks,forget plot] coordinates {(0,0.7854) (1,0.82) (2,0.8533) (3,0.8832) (4,0.8996) (5,0.9156) (6,0.9231)};
\addplot[black,dashed,no marks,forget plot] coordinates {(0,0.95) (6,0.95)};
\end{axis}
\end{tikzpicture}
\hfill\begin{tikzpicture}
\begin{axis}[width=6.6cm,height=5cm,xlabel={Sample size},ylabel={Slope coverage},title={HC3},xtick={0,1,2,3,4,5,6},xticklabels={{25},{50},{100},{200},{400},{800},{1600}}]
% Interval meaning: Wilson 95% interval for coverage over 10000 trials
\addplot+[no marks,forget plot] coordinates {(0,0.8617346296685646) (0,0.8749823635322467)};
\addplot+[no marks,forget plot] coordinates {(1,0.8719508884456331) (1,0.8847584247877647)};
\addplot+[no marks,forget plot] coordinates {(2,0.8874046014402782) (2,0.899493115042325)};
\addplot+[no marks,forget plot] coordinates {(3,0.9026910992750709) (3,0.9139951740560862)};
\addplot+[no marks,forget plot] coordinates {(4,0.9128724734222167) (4,0.923606196741201)};
\addplot+[no marks,forget plot] coordinates {(5,0.9237998614131585) (5,0.9338706679848011)};
\addplot+[no marks,forget plot] coordinates {(6,0.9257617947322282) (6,0.9357072754719226)};
\addplot+[only marks,forget plot] coordinates {(0,0.8685) (1,0.8785) (2,0.8936) (3,0.9085) (4,0.9184) (5,0.929) (6,0.9309)};
\addplot[black,dashed,no marks,forget plot] coordinates {(0,0.95) (6,0.95)};
\end{axis}
\end{tikzpicture}

  \caption{Coverage of nominal 95\% slope intervals on one retained fit per
    trial. Wilson 95\% intervals for the coverage proportion are drawn but are
    mostly narrower than the markers. The dashed line marks the nominal
    level. Each panel is a \texttt{PlotSpec} rendered
  to PGFplots; the manuscript sets a common vertical range.}
  \label{fig:coverage}
\end{figure}

\Cref{fig:coverage} reports the results with Wilson intervals
\citep{wilson-1927}. At $n = 25$, coverage rises from
\jcnum{coverage-hc0-25} for HC0 to \jcnum{coverage-hc3-25} for HC3. This is
consistent with the small-sample advantage of HC3 reported by
\citet{long-ervin-2000}. No estimator reaches the nominal level in this design:
at $n = 1600$ coverage is \jcnum{coverage-hc0-1600} for HC0 and
\jcnum{coverage-hc3-1600} for HC3, because heavy-tailed leverage slows
convergence. A numerically exact covariance computation can therefore
accompany a miscalibrated inferential procedure, which is why STATADIC records
these as separate claims. The study illustrates that distinction; it does not
qualify any HC estimator.

\subsection{Retained qualification records}

The repository also retains prespecified statistical studies run before this
revision. Each uses 10{,}000 trials per scenario, failure-inclusive accounting
and fixed acceptance criteria; for nominal 95\% intervals, the one-sided 95\%
lower bound on coverage must be at least 0.94. \Cref{tab:aft} reproduces the
ordinary AFT records verbatim from the repository's validation guide. Of the
eight cells, \jcvalue{aft-qualified} qualify and \jcvalue{aft-inconclusive} are
inconclusive, and the inconclusive assertions remain failing rather than being
relaxed.

\begin{table}[t]
  \centering
  \small
  \caption{Retained ordinary AFT qualification records: Wald coverage,
  10{,}000 trials per cell, float64.}
  \label{tab:aft}
  % AUTO-GENERATED from docs/validation.md by analysis/generate.py; do not edit.
\begin{tabular}{@{}lrlrrl@{}}
\toprule
Family & $n$ & Entry & Valid & Coverage lower bound & Disposition \\
\midrule
Weibull & 128 & Zero & 10{,}000/10{,}000 & 0.941000 & Qualified \\
Lognormal & 128 & Zero & 10{,}000/10{,}000 & 0.937081 & Inconclusive \\
Weibull & 512 & Zero & 10{,}000/10{,}000 & 0.941722 & Qualified \\
Lognormal & 512 & Zero & 10{,}000/10{,}000 & 0.944717 & Qualified \\
Weibull & 128 & Positive & 9{,}999/10{,}000 & 0.939556 & Inconclusive \\
Lognormal & 128 & Positive & 10{,}000/10{,}000 & 0.940381 & Qualified \\
Weibull & 512 & Positive & 10{,}000/10{,}000 & 0.945853 & Qualified \\
Lognormal & 512 & Positive & 10{,}000/10{,}000 & 0.943684 & Qualified \\
\bottomrule
\end{tabular}

\end{table}

Other records are adverse:
\begin{itemize}
  \item The public Sharpe-ratio interval at $T = 128$ covered
    \jcvalue{sharpe-gaussian}, \jcvalue{sharpe-skewed} and
    \jcvalue{sharpe-dependent} of 10{,}000 Gaussian, skewed and
    dependent-return trials, each reported as excess error.
  \item Spatial survival likelihoods under high censoring miss their numerical
    integration target.
\end{itemize}
These records were not re-executed for this paper, and qualification of a
frozen release remains outstanding.

\subsection{Computation}\label{sec:performance}

\Cref{tab:timing} times STATADIC's linear-kernel V-statistic MMD against an
explicit evaluation of the left-hand side of \cref{thm:linear-mmd} through
three Gram matrices. Both samples have $d = \jcvalue{timing-dimension}$
dimensions and uniform masses, so the equal-mass hypothesis holds, and the two
routes agree to rounding.

The mean route grows linearly with $n$, while the Gram route grows
quadratically: once compiled, it is about \jcvalue{timing-ratio-8000} times
slower at $n = 8000$. First calls include tracing and compilation. Timings are
medians of \jcvalue{timing-repeats} synchronized warm calls on one
\jcvalue{cpu-model} with \jcvalue{cpu-count} logical CPUs. They are
single-machine measurements, not the repository's multi-process benchmark
protocol, and they imply nothing about accelerators, which have not been
measured.

\begin{table}[t]
  \centering
  \small
  \caption{Linear-kernel $\mathrm{MMD}^2$: STATADIC's mean route against explicit
  Gram evaluation, float64 on CPU.}
  \label{tab:timing}
  % AUTO-GENERATED by analysis/generate.py; machine-dependent; do not edit.
\begin{tabular}{@{}rrrrrr@{}}
\toprule
 & \multicolumn{2}{c}{Mean route (ms)} & \multicolumn{2}{c}{Gram route (ms)} & \\
\cmidrule(lr){2-3}\cmidrule(lr){4-5}
$n$ per sample & First call & Warm & First call & Warm & $|\Delta \mathrm{MMD}^2|$ \\
\midrule
500 & 62.1 & 0.02 & 21.6 & 0.13 & \num{1.1e-16} \\
1{,}000 & 55.5 & 0.03 & 21.8 & 0.23 & \num{0} \\
2{,}000 & 54.2 & 0.06 & 23.6 & 0.90 & \num{0} \\
4{,}000 & 59.7 & 0.08 & 53.3 & 32.91 & \num{1.7e-16} \\
8{,}000 & 55.6 & 0.11 & 152.4 & 128.69 & \num{3.9e-16} \\
\bottomrule
\end{tabular}

\end{table}
